% =====================================================================
% DenseNet-121 method, training setup and comparison  (author: Ajit Yadav, Stage 2)
% Every number below comes from results/ in the repository:
%   results/stage2_metrics.csv, results/stage2_comparison.csv, results/stage1_metrics.csv,
%   results/stage2_phaseA/settings.json + val_metrics.csv, results/stage2_phaseB/tuning.csv,
%   results/stage2_densenet/settings.json
% Figures to upload to Overleaf (keep these paths):
%   figures/stage2/confusion_matrices.png   <- results/stage2_densenet/confusion_matrices.png
%   figures/stage2/roc.png                  <- results/stage2_densenet/roc.png
% Labels used from Abhishek's section: sec:data-splits, sec:data-preproc, sec:baseline, tab:baseline, sec:data-limits
% =====================================================================
\section{DenseNet-121 Transfer Learning}
\label{sec:densenet}

\subsection{Method}
\label{sec:densenet-method}
The stronger model is DenseNet-121~\cite{huang2017densenet} initialised with ImageNet~\cite{deng2009imagenet} weights, following CheXNet~\cite{rajpurkar2017chexnet}.
We replace the original 1000-class layer by a dropout layer (rate 0.2) and one linear layer with a single output logit, where the sigmoid of the logit is the probability of pneumonia.
The network has 6.95\,M parameters.
The data pipeline is exactly the one used for the baseline (\cref{sec:data-preproc}): grayscale images copied to three channels, plain resize to $224\times224$,
ImageNet mean and standard deviation, and a train-only rotation of $\pm10^\circ$ without horizontal flip.
It also uses the same patient-level splits (\cref{sec:data-splits}), the same weighted loss, and the same metrics as the baseline, so the two models differ only in the architecture and its initialisation.

\subsection{Training setup}
\label{sec:densenet-training}
We train in two phases, as in CheXNet-style transfer learning.
In \emph{Phase~A}, the backbone is frozen (its BatchNorm layers are kept in evaluation mode, so their statistics do not change) and only the new output layer, 1,025 parameters, is trained
with Adam (learning rate $10^{-3}$, weight decay $10^{-4}$, batch size 32) for up to 20 epochs.
The best epoch (19) has a validation loss of 0.065 and a validation AUC of 0.990.
In \emph{Phase~B}, we start from the Phase~A weights, unfreeze the whole network and fine-tune all 6.95\,M parameters with a low learning rate
(Adam, weight decay $10^{-4}$, for at most 15 epochs).
Both phases minimise \texttt{BCEWithLogitsLoss} with the class weight computed on the training split only ($\text{pos\_weight}=0.347$),
halve the learning rate when the validation loss stops improving (patience 2), and stop early on the validation loss (patience 5 in Phase~A, 4 in Phase~B).
Mixed precision is used on the GPU, and all runs use seed 42.

\paragraph{Tuning.}
We tuned the learning rate and the batch size of Phase~B with three runs (\cref{tab:tuning}).
The model is selected by the lowest \emph{validation} loss; the test set is not used for any choice and is evaluated once, for the selected model only.
The selected run uses a learning rate of $10^{-4}$ and batch size 16 (best epoch 3).

\begin{table}[t]
  \centering
  \small
  \begin{tabular}{lrrrr}
    \toprule
    Run (Phase B) & LR & Batch & Val loss & Val AUC \\
    \midrule
    \textbf{lr1e-4\_bs16} & $10^{-4}$ & 16 & \textbf{0.0139} & 0.9995 \\
    lr1e-4                & $10^{-4}$ & 32 & 0.0206 & 0.9994 \\
    lr5e-5                & $5\cdot10^{-5}$ & 32 & 0.0278 & 0.9984 \\
    \bottomrule
  \end{tabular}
  \caption{Phase~B tuning on the validation set (patient-level split). The run with the lowest validation loss (bold) is the final model.}
  \label{tab:tuning}
\end{table}

\subsection{Results and comparison with the baseline}
\label{sec:densenet-results}
\Cref{tab:comparison} compares the final DenseNet-121 with the baseline CNN (\cref{sec:baseline}) on the same official test set with the same metrics (threshold 0.5).
DenseNet-121 raises accuracy from 0.801 to 0.873, F1 from 0.861 to 0.908 and AUC from 0.902 to 0.971.
The number of normal images predicted as pneumonia falls from 117 to 78 of 234, and only one of 390 pneumonia images is missed (baseline: 7).
The pretrained backbone provides better features than a network trained from scratch on about 4,500 images, which is the usual benefit of transfer learning.
\Cref{fig:densenet-cm} shows the confusion matrices and \cref{fig:densenet-roc} the ROC curves.

\begin{table}[t]
  \centering
  \small
  \begin{tabular}{lrrrrrrr}
    \toprule
    Model & Acc. & Prec. & Rec. & F1 & AUC & FP & FN \\
    \midrule
    Custom CNN   & 0.801 & 0.766 & 0.982 & 0.861 & 0.902 & 117 & 7 \\
    DenseNet-121 & \textbf{0.873} & \textbf{0.833} & \textbf{0.997} & \textbf{0.908} & \textbf{0.971} & \textbf{78} & \textbf{1} \\
    \bottomrule
  \end{tabular}
  \caption{Test-set results of the baseline CNN and DenseNet-121 (unmasked images, threshold 0.5; 234 normal and 390 pneumonia images). FP: normal images predicted as pneumonia; FN: missed pneumonia images.}
  \label{tab:comparison}
\end{table}

\begin{figure}[t]
  \centering
  \includegraphics[width=\linewidth]{figures/stage2/confusion_matrices.png}
  \caption{Confusion matrices of DenseNet-121 on the validation and test sets (threshold 0.5).}
  \label{fig:densenet-cm}
\end{figure}

\begin{figure}[t]
  \centering
  \includegraphics[width=0.7\linewidth]{figures/stage2/roc.png}
  \caption{ROC curves of DenseNet-121 on the validation and test sets.}
  \label{fig:densenet-roc}
\end{figure}

\paragraph{Validation--test gap.}
Like the baseline, DenseNet-121 performs much better on the validation set (accuracy 0.985, AUC 0.9995) than on the test set (accuracy 0.873, AUC 0.971).
The validation set is drawn from the same distribution as the training data, whereas the official test set differs from it, so the test metrics are the ones that matter.
DenseNet-121 narrows the gap (0.111 accuracy difference, baseline: 0.167) but does not remove it, and 78 normal test images are still predicted as pneumonia.
Whether these errors come from image regions outside the lungs is examined with Grad-CAM in the next section. % TODO: link to Aman's Grad-CAM section (Stage 3)

\paragraph{Limitations.}
All images come from one medical centre (\cref{sec:data-limits}), the three tuning runs cover only a small part of the hyperparameter space,
and all numbers come from a single training run with seed 42, so small differences between models should not be over-interpreted.
% TODO (Himanshu, Stage 4): add the lung-masked DenseNet-121 row to \cref{tab:comparison} and the attention results.
